\documentclass{jsarticle}
\usepackage{amssymb}
\usepackage{amsthm}
\usepackage{amsmath}
\newtheorem{theorem}{定理}
\newtheorem{proposition}[theorem]{命題}
\newtheorem{corollary}[theorem]{系}
\newtheorem{lemma}[theorem]{補題}
\newtheorem{defnition}[theorem]{定義}
\newtheorem{example}[theorem]{例}
\newtheorem{fact}[theorem]{事実}
\newtheorem*{claim}{主張}
\renewcommand{\proofname}{証明}
\title{ウリゾーンの距離化可能定理}
\author{電波通信}
\date{\today}
\begin{document}
\maketitle
この文章ではウリゾーンの距離化可能定理を証明する。ギリシャ文字の$\phi$と空集合を表す$\O$を間違えないように注意せよ。

\begin{defnition}
位相空間$X$が正規であるとは交わらない二つの閉集合の任意の組$F,G$に対して開集合$U,V$が存在して
\[
F\subset U,\ G\subset V,\ U\cap V =\emptyset
\]を充たすことである。
\end{defnition}
次が成立する。

\begin{fact}[ウリゾーンの補題]
$(X,\mathfrak{O})$を正規空間とし、$F,G$をその互いに素な２つの閉集合とすると連続関数$f:X\rightarrow [0,1]$が存在し
\[
f|F=1,\ f|G=0
\]つまり互いに素な２つの閉集合は関数で分離できる。
\end{fact}
ウリゾーンの補題は次のように変形して使われることも多い。

\begin{fact}[ウリゾーンの補題II]
$(X,\mathfrak{O})$を正規空間とし、$A$を閉集合、$O$を開集合とし$A\subset U$とすると連続関数$f:X\rightarrow [0,1]$が存在し
\[
f|A=0,\ f|X\setminus O=1
\]を充たす
\end{fact}これの証明は明らかである。

\begin{defnition}[始位相]
$X$を集合$\{(Y_{\lambda},\mathfrak{O}_{\lambda})\}_{\lambda \in \Lambda}$を位相空間の族とし、各$\lambda$ 毎に写像
\[
f_{\lambda}:X\rightarrow Y_{\lambda}
\]が与えられているとするこの時$X$に写像族$F=\{f_{\lambda}\}_{\lambda \in \Lambda}$をすべて連続にするような最弱の位相が存在する。それを$F=\{f_{\lambda}\}_{\lambda \in \Lambda}$から誘導された始位相と呼ぶ。此処では記号$\mathfrak{I}(F)$で表すことにする。この位相は準開基$\mathfrak{M}=\{f_{\lambda}^{-1}(U)\ |\ U\in \mathfrak{O}_{\lambda},\ \lambda\in \Lambda \}$により生成されている
\end{defnition}$\mathfrak{I}$は$I$のフラクトゥールである。

\begin{fact}[始位相の埋め込み定理]\label{main}位相空間$(X,\mathfrak{O})$には$\{f_{\lambda}\}_{\lambda \in \Lambda}$から定まる始位相が入っているとしこの関数族の終域を$\{(Y_{\lambda},\mathfrak{O}_{\lambda})\}_{\lambda \in \Lambda}$とするこの時関数
\[
\phi=\prod_{\lambda\in \Lambda}f_{\lambda}:X\rightarrow \prod_{\lambda\in \Lambda}Y_{\lambda \in \Lambda}
\]について$\mathfrak{I}(F)=\mathfrak{I}(\phi)$が成立する。
また$\phi$が単射ならばこれは埋め込みになっているつまり$(X,\mathfrak{O})$は$\prod_{\lambda\in \Lambda}Y_{\lambda \in \Lambda}$の部分空間$\phi(X)$と同相となる。
\end{fact}

\begin{fact}[距離空間の可算直積]距離空間の可算直積は距離空間となる。
\end{fact}

\begin{defnition}[ヒルベルトキューブ]
単位区間$I=[0,1]$の可算直積$I^{\aleph_0}$をヒルベルトキューブという。上の事実からこの空間は距離空間である。
\end{defnition}

次の事実は当たり前である。
\begin{fact}
距離空間の部分空間は距離空間である。
\end{fact}

\begin{defnition}[第二可算公理]位相空間が高々可算な開基を持つとき、つまり
\[
\text{card$(\mathfrak{B})$}\le \aleph_0
\]となる開基$\mathfrak{B}$を持つとき、その空間は第二可算公理を充たすという。
\end{defnition}

\begin{theorem}[ウリゾーンの距離化可能定理]第二可算公理を充たす$T_4$空間は距離化可能、つまりその空間の位相と両立する距離を持つ。
\end{theorem}
ウリゾーンの補題と始位相の埋め込み定理を使い空間をヒルベルトキューブに埋め込むのが証明のアイディアである。

\begin{proof}
$X$が第二可算性から存在がわかる高々可算な開基のひとつを$\mathfrak{B}$としよう。
この時$\mathfrak{B}\times \mathfrak{B}$も高々可算である。この集合の部分集合$\mathfrak{C}$を
\[
\mathfrak{C}=\{(P,Q)\in \mathfrak{B}\times\mathfrak{B}\ |\ \bar{P}\subset Q\}
\]と定義する。もちろん$\mathfrak{C}$も高々可算である。
次に（ウリゾーンの補題II）を使って各$(P,Q)\in \mathfrak{C}$に対して
\[
f|\bar{P}=0,\ f|X\setminus Q=1
\]
となる連続関数$f$が存在する。ところで$\mathfrak{C}$は高々可算であるから自然数$\mathbb{N}$を用いて
\[
\mathfrak{C}=\{C_n\}_{n\in\mathbb{N}}
\]と書ける。$C_n$の具体的な中身が見たい時には、$C_n=(P_n,Q_n)$と書くことにしよう。
各$C_n=(P_n,Q_n)$に対してウリゾーンの補題IIから存在が保証される上のような関数ののひとつを$f_n$と書こう。さて始位相の埋め込み定理を利用するために次を示そう。
\begin{claim}
\[
\mathfrak{O}=\mathfrak{I}(\{f_n\}_{n\in \mathbb{N}})
\]
つまり$X$の位相は$\{f_n\}_{n\in \mathbb{N}}$から誘導される始位相と一致するのである。
\end{claim}
\begin{proof}（主張の証明）各$f_n$は連続なので明らかに
\[
\mathfrak{O}\supset \mathfrak{I}(\{f_n\}_{n\in \mathbb{N}})
\]逆向きの包含関係を示そう。今$B\in \mathfrak{B}$を任意に取ると$X$の正規性と$\mathfrak{B}$が開基であることから
\[
\bar{A}\subset B
\]
となる$A\in \mathfrak{B}$が存在する。よって$(A,B)\in \mathfrak{C}$なのである自然数$n$が存在して$C_n=(A,B)$と書ける。さて$f_n$は定義から
\[
f_n|\bar{A}=0,\ f|X\setminus B=1
\]となる。よって
\[
\bar{A}\subset f_n^{-1}([0,\frac{1}{2}))\subset B
\]
が成り立ち、始位相の準開基の形と$B\in\mathfrak{B}$の任意性と$\mathfrak{B}$が開基であることから
\[
\mathfrak{O}\subset \mathfrak{I}(\{f_n\}_{n\in \mathbb{N}})
\]
[主張の証明終]\end{proof} 
[定理の証明の続]
さて次に$F=\prod_{n}f_n$が単射であることを示そう。$x\neq y$とすると、$T_1$から$B\in \mathfrak{B}$が存在して$x\in B, \ y\not\in B$を充たす。先と同様に正規性と$\mathfrak{B}$が開基であることから
$\bar{A}\subset B,\ A\in \mathfrak{B}$となる$A$が存在し、ある自然数$n$が存在して
$C_n=(A,B)\in \mathfrak{C}$となる。この時$f_n$は
\[
f_n|\bar{A}=0,\ f|X\setminus B=1
\]をみたし、特に
\[
f(x)=0,\ f(y)=1
\]
となる。よって
\[
x\neq y\Rightarrow F(x)\neq F(y)
\]よって$F$は単射。\\ 
以上から始位相の埋め込み定理を使うことができ、$X$は$[0,1]^{\aleph_0}$に埋め込まれる。つまり、ヒルベルトキューブの部分空間と同相である。よって明らかに$X$は距離化可能である。
\end{proof}

\begin{corollary}
コンパクトハウスドルフ空間に於いて距離化可能であることと第二可算公理を充たすことは同値である。
\end{corollary}

\begin{proof}
コンパクトハウスドルフ空間は正則であることとコンパクト距離空間が常に第二可算公理を充たすことを踏まえればすぐにわかる。
\end{proof}

\begin{fact}
第二可算公理を充たす正則空間は正規である。
\end{fact}

\begin{corollary}
\footnote{より強く、第二可算公理を充たす局所コンパクトハウスドルフ空間は完備距離付け可能であることが知られている。完備距離空間の開集合に同相だからであるが、それはまた別の話。}
局所コンパクトハウスドルフ空間は第二可算公理を満たせば距離化可能である。
\end{corollary}

\begin{proof}局所コンパクトハウスドルフ空間は完全正則だから先の事実とウリゾーンの距離化可能定理からすぐにわかる。
\end{proof}



\begin{thebibliography}{9}
\bibitem{1}森田紀一,位相空間論,岩波全書331,岩波書店1981
\bibitem{2}彌永昌吉・彌永健一,集合と位相,岩波基礎数学選書,2002
\end{thebibliography}







\end{document}